% ---------------------------------------------------------------------------
% Haast Autonomous - Summer Engineering Intern (aircraft development)
% Tailored from bases/simulation_robotics.tex on 2026-08-03
% Worksheet: notes.md in this folder. Posting: job_description.txt
%
% Positioning: they name PX4 and Formula SAE in their own nice-to-haves and you
% have both. Argus PX4 work leads. Skills regrouped so flight/autonomy is the
% first thing read. SR-Wireless-CAN swapped in for Posture Detection - hardware
% bring-up under ambiguity is the trait they say they hire for.
%
% Do NOT put the company name anywhere in the rendered document. It belongs in
% the filename only.
% ---------------------------------------------------------------------------
\documentclass[10pt,letterpaper]{article}
\usepackage{resume}

\begin{document}
\bodyfont

\resumeheader
\educationsection

\sectionheader{TECHNICAL SKILLS}

\techline{\textbf{Flight \& Autonomy}: PX4 flight stack, flight-control firmware, visual-inertial odometry, IMU sensor fusion, state estimation, flight testing, telemetry log analysis}
\techline{\textbf{Controls}: PID and fixed-point PI, gain tuning, anti-windup, saturation, slew-rate limiting, deterministic real-time execution}
\techline{\textbf{Simulation \& Test}: Physics-based vehicle simulation, software-in-the-loop replay, regression suites, parameter sweeps, headless batch runs, pytest}
\techline{\textbf{Languages \& Tools}: Python, C/C++, NumPy, pandas, Matplotlib, OpenCV, CAN/DBC, GDB, BusMaster, MoTeC i2 Pro, Git}

\sectionheader{EXPERIENCE}

\roleline{Software Engineer Intern}{Feb 2026 - Present}
\orgline{Argus Defense \textbar{} San Jose, CA}
\begin{itemize}
  \item Developed and validated embedded flight-control firmware in C on PX4, iterating control gains against logged flight telemetry
  \item Integrated and tuned visual-inertial odometry for onboard state estimation and evaluated performance against logged telemetry
  \item Tuned PX4 PID flight-control loops to reduce vertical and horizontal drift during autonomous position hold, iterating through repeatable flight tests and log review
\end{itemize}

\entryspace

\roleline{Software Lead}{Jun 2026 - Present}
\orgline{\spartanracing}
\begin{itemize}
  \item Lead validation strategy for vehicle controls, telemetry, diagnostics, simulation, and on-vehicle test workflows
  \item Define testable software requirements and coordinate reviews across controls, dashboard, simulation, and software-in-the-loop projects
\end{itemize}

\entryspace

\roleline{Software Controls Engineer}{Jul 2025 - Jun 2026}
\orgline{\spartanracing}
\begin{itemize}
  \item Implemented a torque-to-power controller with fixed-point PI control, saturation, anti-windup, and slew limits for deterministic real-time execution
  \item Built Python telemetry parsers and a software-in-the-loop replay harness to validate controller changes against recorded logs before hardware testing
  \item Planned and executed on-vehicle tests, measured response and latency, and turned results into parameter changes and regression checks
\end{itemize}

\entryspace

\roleline{Software Intern}{Aug 2024 - Jun 2025}
\orgline{\spartanracing}
\begin{itemize}
  \item Extracted a desktop simulator and regression suite to reproduce field scenarios, run controller sweeps, and prevent regressions
  \item \textbf{Results:} 1st in endurance at MIS EV 2025
\end{itemize}

\sectionheader{PROJECT EXPERIENCE}

\roleline{SR-Wireless-CAN --- Wireless Firmware Flashing \& Telemetry}{2026}
\orgline{\spartanracing \textbar{} Backend lead}
\techline{Python, python-can, bus trace analysis, FastAPI, WebSockets, SQLite, Docker}
\begin{itemize}
  \item Proposed and implemented wireless firmware flashing for the vehicle control unit, decoding the undocumented flashing sequence from bus trace captures and reimplementing it in Python
  \item Led backend development for the telemetry platform, streaming live signals to a browser dashboard with recorded-trace replay for offline test review
\end{itemize}

\entryspace

\roleline{Vehicle Dynamics Simulation}{Jul 2025}
\orgline{\spartanracing}
\techline{Python, Pygame, Matplotlib, vehicle dynamics, PID control, CSV logging}
\begin{itemize}
  \item Built a physics-based simulator for torque, drag, acceleration, and steering to evaluate control changes before committing them to hardware
  \item Added headless batch runs, seeded scenarios, CSV summaries, and plots for repeatable offline comparison between builds
\end{itemize}

\end{document}
